\documentclass[11pt]{article}
\usepackage[T1]{fontenc}
\usepackage{amsmath,amssymb,amsthm}
\usepackage[margin=1in]{geometry}
\usepackage{booktabs}
\usepackage{hyperref}
\hypersetup{
  pdftitle={A Certified Prime Factor of a22 in a Quadratic Recurrence},
  pdfauthor={Xiaojian Zeng},
  pdfsubject={A deterministic primality and modular-divisibility certificate},
  pdfkeywords={integer recurrence, prime factor, Pocklington certificate}
}

\newtheorem{theorem}{Theorem}
\newtheorem{lemma}[theorem]{Lemma}

\title{A Certified Prime Factor of $a_{22}$ in the Recurrence
  $a_{n+1}=a_n^2+2$}
\author{Xiaojian Zeng\\
  PBC School of Finance, Tsinghua University, 100190\\
  \texttt{xiaojianzeng@foxmail.com}}
\date{}

\begin{document}
\maketitle

\begin{abstract}
For the integer sequence $a_0=1$ and $a_{n+1}=a_n^2+2$, we give the
explicit prime
\[
q=108099707
\]
and exact certificates showing that $q\mid a_{22}$ and $1<q<a_{22}$.
Consequently $a_{22}$ is composite.  No claim is made about the remaining
cofactor, later terms of the sequence, or novelty or priority.
\end{abstract}

\section{Statement and context}

\begin{theorem}\label{thm:main}
Let $a_0=1$ and $a_{n+1}=a_n^2+2$ for $n\geq0$.  Then
$q=108099707$ is prime, $q\mid a_{22}$, and $1<q<a_{22}$.  In particular,
$a_{22}$ is composite.
\end{theorem}

A 2026 preprint by Zeng concerns the same recurrence and, according to its
publisher-deposited abstract, reports the factor $7795379$ of $a_{20}$ and
makes no assertion for $n\geq22$ \cite{Zeng2026SSRN7395382}.  Its full text
was not available in the bounded source search for this note, and no result
from it is used below.  The precise delta from that available metadata is the
certified factor of $a_{22}$ in Theorem~\ref{thm:main}.  The source search was
not exhaustive, so this comparison is not a claim of novelty or priority.

For broader arithmetic-dynamical context, Jones proved that if
$f(x)=x^2+k$ with $k\in\mathbb Z\setminus\{-1\}$, then the set of primes
dividing a nonzero term in any integral orbit of $f$ has natural density zero
\cite[Theorem~1.2]{Jones2008Density}.  Applied here with $k=2$ and initial
point $1$, that asymptotic result does not identify a divisor of $a_{22}$ and
is not used in the proof.

We use the following complete-factorization criterion.  Brillhart, Lehmer,
and Selfridge state the corresponding $N-1$ criterion with the weaker
condition
$b_p^{(N-1)/p}\not\equiv1\pmod N$
\cite[Theorem~1]{BrillhartLehmerSelfridge1975}; the greatest-common-divisor
condition below is stronger.  A proof is included to keep the certificate
logically self-contained.

\begin{lemma}[Full-factorization criterion]\label{lem:criterion}
Let $N>1$, and suppose the complete prime factorization of
$N-1=\prod_p p^{e_p}$ is known.  If, for every prime $p\mid N-1$, there is an
integer $b_p$ such that
\[
b_p^{N-1}\equiv1\pmod N,
\qquad
\gcd\bigl(b_p^{(N-1)/p}-1,N\bigr)=1,
\]
then $N$ is prime.
\end{lemma}

\begin{proof}
Let $\ell$ be any prime divisor of $N$.  The displayed conditions imply that
the multiplicative order of $b_p$ modulo $\ell$ divides $N-1$ but does not
divide $(N-1)/p$.  Hence its $p$-adic valuation is $e_p$, so
$p^{e_p}\mid\ell-1$.  This holds for every $p\mid N-1$, giving
$N-1\mid\ell-1$.  Since $\ell\leq N$, it follows that $\ell=N$, and therefore
$N$ is prime.
\end{proof}

\section{Primality certificate}

First,
\[
34361-1=34360=2^3\cdot5\cdot859.
\]
The integer $859$ is not divisible by any prime at most
$\lfloor\sqrt{859}\rfloor=29$: for the primes
\[
2,3,5,7,11,13,17,19,23,29
\]
its respective nonzero residues are
\[
1,1,4,5,1,1,9,4,8,18.
\]
Thus $859$ is prime.  The following rows give the exact data required by
Lemma~\ref{lem:criterion}; the third and fourth columns are reduced modulo
$N$.

\begin{center}
\begin{tabular}{rrrrr}
\toprule
$N$ & $p$ & $b_p^{N-1}$ & $b_p^{(N-1)/p}$ &
$\gcd(b_p^{(N-1)/p}-1,N)$ \\
\midrule
$34361$ & $2$   & $1$ & $34360$ & $1$ \\
$34361$ & $5$   & $1$ & $15387$ & $1$ \\
$34361$ & $859$ & $1$ & $1951$  & $1$ \\
\bottomrule
\end{tabular}
\end{center}
Here $b_2=3$ and $b_5=b_{859}=2$.  Lemma~\ref{lem:criterion} therefore
proves that $34361$ is prime.  Next,
\[
q-1=108099706=2\cdot11^2\cdot13\cdot34361.
\]
Using $b_p=2$ in every row gives

\begin{center}
\begin{tabular}{rrrrr}
\toprule
$N$ & $p$ & $b_p^{N-1}$ & $b_p^{(N-1)/p}$ &
$\gcd(b_p^{(N-1)/p}-1,N)$ \\
\midrule
$108099707$ & $2$     & $1$ & $108099706$ & $1$ \\
$108099707$ & $11$    & $1$ & $54894374$  & $1$ \\
$108099707$ & $13$    & $1$ & $39708810$  & $1$ \\
$108099707$ & $34361$ & $1$ & $99912603$  & $1$ \\
\bottomrule
\end{tabular}
\end{center}
The elementary primes $2,11,13$, together with the preceding certification of
$34361$, give the complete prime factorization of $q-1$.  A second application
of Lemma~\ref{lem:criterion} proves that $q$ is prime.

\section{Divisibility and properness}

Starting with $r_0=1$, compute exactly
\[
r_{n+1}=(r_n^2+2)\bmod q,
\qquad 0\leq n<22,
\]
where every displayed residue is the least nonnegative one.  The full chain is

\begin{center}
\begin{tabular}{rr@{\qquad}rr}
\toprule
$n$ & $r_n$ & $n$ & $r_n$ \\
\midrule
$0$  & $1$         & $12$ & $13474761$  \\
$1$  & $3$         & $13$ & $51643108$  \\
$2$  & $11$        & $14$ & $36226811$  \\
$3$  & $123$       & $15$ & $44888898$  \\
$4$  & $15131$     & $16$ & $33268166$  \\
$5$  & $12747749$  & $17$ & $18262376$  \\
$6$  & $4130680$   & $18$ & $80442749$  \\
$7$  & $59509522$  & $19$ & $68864772$  \\
$8$  & $53448106$  & $20$ & $83683223$  \\
$9$  & $101159942$ & $21$ & $103702901$ \\
$10$ & $81091708$  & $22$ & $0$         \\
$11$ & $50080613$  &      &             \\
\bottomrule
\end{tabular}
\end{center}
By induction on $n$, the defining recurrence gives
$a_n\equiv r_n\pmod q$ for $0\leq n\leq22$.  Since $r_{22}=0$, we have
$q\mid a_{22}$.

Direct integer recursion gives
\[
(a_0,a_1,a_2,a_3,a_4,a_5)
=(1,3,11,123,15131,228947163).
\]
Thus $1<q<a_5$.  Every $a_n$ is positive, and
$a_{n+1}=a_n^2+2>a_n$, so $a_5<a_{22}$.  This proves
$1<q<a_{22}$ and completes the proof of Theorem~\ref{thm:main}.

\section*{Reproducibility and assistance disclosure}

The supplementary standard-library Python program
\texttt{check\_certificate.py} recomputes both complete factorizations, every
modular power and greatest common divisor in the primality certificate, all 22
modular-recurrence transitions, and the exact value of $a_5$.  It also performs
an independent exhaustive divisibility test for $q$ by every integer from $2$
through $\lfloor\sqrt q\rfloor=10397$.

The general two-level pattern of generating Pocklington witnesses externally
and checking them in a smaller proof environment has also been implemented in
the Coq proof assistant \cite{GregoireTheryWerner2006}.  The checker supplied
with this note is an ordinary Python program using exact integer arithmetic;
no Coq formalization is claimed here.

The factor was found by a computer search over candidate prime divisors.
Computer algebra was not used in the proof: the released checker uses exact
integer arithmetic only.  An AI language model assisted with search-code
development, certificate organization, and manuscript drafting.  The numerical
certificates were then checked by two deterministic programs and independently
recomputed during mathematical review.  The theorem rests on the exact
certificates printed above, not on probabilistic primality testing, floating
point computation, or the literature search.

\bibliographystyle{plain}
\bibliography{references}

\end{document}
